\section{AHSME 1954}
        \begin{enumerate}
        	\item (\textit{AHSME 1954}) The square of $5-\sqrt{y^2-25}$ is: 


 $\textbf{(A)}\ y^2-5\sqrt{y^2-25} \qquad \textbf{(B)}\ -y^2 \qquad \textbf{(C)}\ y^2 \\ \textbf{(D)}\ (5-y)^2\qquad\textbf{(E)}\ y^2-10\sqrt{y^2-25}$


	\item (\textit{AHSME 1954}) The equation $\frac{2x^2}{x-1}-\frac{2x+7}{3}+\frac{4-6x}{x-1}+1=0$ can be transformed by eliminating fractions to the equation $x^2-5x+4=0$. 
The roots of the latter equation are $4$ and $1$. Then the roots of the first equation are: 


 $\textbf{(A)}\ 4 \text{ and }1 \qquad \textbf{(B)}\ \text{only }1 \qquad \textbf{(C)}\ \text{only }4 \qquad \textbf{(D)}\ \text{neither 4 nor 1}\qquad\textbf{(E)}\ \text{4 and some other root}$


	\item (\textit{AHSME 1954}) If $x$ varies as the cube of $y$, and $y$ varies as the fifth root of $z$, then $x$ varies as the nth power of $z$, where n is:


 $\textbf{(A)}\ \frac{1}{15} \qquad\textbf{(B)}\ \frac{5}{3} \qquad\textbf{(C)}\ \frac{3}{5} \qquad\textbf{(D)}\ 15\qquad\textbf{(E)}\ 8$


	\item (\textit{AHSME 1954}) If the Highest Common Divisor of $6432$ and $132$ is diminished by $8$, it will equal: 


 $\textbf{(A)}\ -6 \qquad \textbf{(B)}\ 6 \qquad \textbf{(C)}\ -2 \qquad \textbf{(D)}\ 3 \qquad \textbf{(E)}\ 4$


	\item (\textit{AHSME 1954}) A regular hexagon is inscribed in a circle of radius $10$ inches. Its area is: 


 $\textbf{(A)}\ 150\sqrt{3} \text{ sq. in.} \qquad \textbf{(B)}\ \text{150 sq. in.} \qquad \textbf{(C)}\ 25\sqrt{3}\text{ sq. in.}\qquad\textbf{(D)}\ \text{600 sq. in.}\qquad\textbf{(E)}\ 300\sqrt{3}\text{ sq. in.}$


	\item (\textit{AHSME 1954}) The value of $\frac{1}{16}a^0+\left (\frac{1}{16a} \right )^0- \left (64^{-\frac{1}{2}} \right )- (-32)^{-\frac{4}{5}}$ is: 


 $\textbf{(A)}\ 1 \frac{13}{16} \qquad \textbf{(B)}\ 1 \frac{3}{16} \qquad \textbf{(C)}\ 1 \qquad \textbf{(D)}\ \frac{7}{8}\qquad\textbf{(E)}\ \frac{1}{16}$


	\item (\textit{AHSME 1954}) A housewife saved $\textdollar{2.50}$ in buying a dress on sale. If she spent $\textdollar{25}$ for the dress, she saved about: 


 $\textbf{(A)}\ 8 \% \qquad \textbf{(B)}\ 9 \% \qquad \textbf{(C)}\ 10 \% \qquad \textbf{(D)}\ 11 \% \qquad \textbf{(E)}\ 12\%$


	\item (\textit{AHSME 1954}) The base of a triangle is twice as long as a side of a square and their areas are the same. 
Then the ratio of the altitude of the triangle to the side of the square is: 


 $\textbf{(A)}\ \frac{1}{4} \qquad \textbf{(B)}\ \frac{1}{2} \qquad \textbf{(C)}\ 1 \qquad \textbf{(D)}\ 2 \qquad \textbf{(E)}\ 4$


	\item (\textit{AHSME 1954}) A point $P$ is outside a circle and is $13$ inches from the center. A secant from $P$ cuts the circle at $Q$ and $R$ 
so that the external segment of the secant $PQ$ is $9$ inches and $QR$ is $7$ inches. The radius of the circle is: 


 $\textbf{(A)}\ 3" \qquad \textbf{(B)}\ 4" \qquad \textbf{(C)}\ 5" \qquad \textbf{(D)}\ 6"\qquad\textbf{(E)}\ 7"$


	\item (\textit{AHSME 1954}) The sum of the numerical coefficients in the expansion of the binomial $(a+b)^6$ is: 


 $\textbf{(A)}\ 32 \qquad \textbf{(B)}\ 16 \qquad \textbf{(C)}\ 64 \qquad \textbf{(D)}\ 48 \qquad \textbf{(E)}\ 7$


	\item (\textit{AHSME 1954}) A merchant placed on display some dresses, each with a marked price. He then posted a sign “$\frac{1}{3}$ off on these dresses.” 
The cost of the dresses was $\frac{3}{4}$ of the price at which he actually sold them. Then the ratio of the cost to the marked price was: 


 $\textbf{(A)}\ \frac{1}{2} \qquad \textbf{(B)}\ \frac{1}{3} \qquad \textbf{(C)}\ \frac{1}{4} \qquad \textbf{(D)}\ \frac{2}{3}\qquad\textbf{(E)}\ \frac{3}{4}$


	\item (\textit{AHSME 1954}) The solution of the equations 


 
 \begin{align*}2x-3y &=7 \\ 4x-6y &=20\end{align*}



 is: 


 $\textbf{(A)}\ x=18, y=12 \qquad  \textbf{(B)}\ x=0, y=0 \qquad  \textbf{(C)}\ \text{There is no solution} \\  \textbf{(D)}\ \text{There are an unlimited number of solutions}\qquad \textbf{(E)}\ x=8, y=5$


	\item (\textit{AHSME 1954}) A quadrilateral is inscribed in a circle. If angles are inscribed in the four arcs cut off 
by the sides of the quadrilateral, their sum will be: 


 $\textbf{(A)}\ 180^\circ \qquad \textbf{(B)}\ 540^\circ \qquad \textbf{(C)}\ 360^\circ \qquad \textbf{(D)}\ 450^\circ\qquad\textbf{(E)}\ 1080^\circ$


	\item (\textit{AHSME 1954}) When simplified $\sqrt{1+ \left (\frac{x^4-1}{2x^2} \right )^2}$ equals: 


 $\textbf{(A)}\ \frac{x^4+2x^2-1}{2x^2} \qquad \textbf{(B)}\ \frac{x^4-1}{2x^2} \qquad \textbf{(C)}\ \frac{\sqrt{x^2+1}}{2}\\ \textbf{(D)}\ \frac{x^2}{\sqrt{2}}\qquad\textbf{(E)}\ \frac{x^2}{2}+\frac{1}{2x^2}$


	\item (\textit{AHSME 1954}) $\log 125$ equals: 


 $\textbf{(A)}\ 100 \log 1.25 \qquad \textbf{(B)}\ 5 \log 3 \qquad \textbf{(C)}\ 3 \log 25  \\ \textbf{(D)}\ 3 - 3\log 2 \qquad \textbf{(E)}\ (\log 25)(\log 5)$


	\item (\textit{AHSME 1954}) If $f(x) = 5x^2 - 2x - 1$, then $f(x + h) - f(x)$ equals: 


 $\textbf{(A)}\ 5h^2 - 2h \qquad \textbf{(B)}\ 10xh - 4x + 2 \qquad \textbf{(C)}\ 10xh - 2x - 2 \\ \textbf{(D)}\ h(10x+5h-2)\qquad\textbf{(E)}\ 3h$


	\item (\textit{AHSME 1954}) The graph of the function $f(x) = 2x^3 - 7$ goes: 


 $\textbf{(A)}\ \text{up to the right and down to the left} \\ \textbf{(B)}\ \text{down to the right and up to the left}\\ \textbf{(C)}\ \text{up to the right and up to the left}\\ \textbf{(D)}\ \text{down to the right and down to the left}\\ \textbf{(E)}\ \text{none of these ways.}$


	\item (\textit{AHSME 1954}) Of the following sets, the one that includes all values of $x$ which will satisfy $2x - 3 > 7 - x$ is: 


 $\textbf{(A)}\ x > 4 \qquad \textbf{(B)}\ x < \frac {10}{3} \qquad \textbf{(C)}\ x = \frac {10}{3} \qquad \textbf{(D)}\ x >\frac{10}{3}\qquad\textbf{(E)}\ x < 0$


	\item (\textit{AHSME 1954}) If the three points of contact of a circle inscribed in a triangle are joined, the angles of the resulting triangle: 


 $\textbf{(A)}\ \text{are always equal to }60^\circ\\ \textbf{(B)}\ \text{are always one obtuse angle and two unequal acute angles}\\ \textbf{(C)}\ \text{are always one obtuse angle and two equal acute angles}\\ \textbf{(D)}\ \text{are always acute angles}\\ \textbf{(E)}\ \text{are always unequal to each other}$


	\item (\textit{AHSME 1954}) The equation $x^3+6x^2+11x+6=0$ has: 


 $\textbf{(A)}\ \text{no negative real roots}\qquad\textbf{(B)}\ \text{no positive real roots}\qquad\textbf{(C)}\ \text{no real roots}\\ \textbf{(D)}\ \text{1 positive and 2 negative roots}\qquad\textbf{(E)}\ \text{2 positive and 1 negative root}$


	\item (\textit{AHSME 1954}) The roots of the equation $2\sqrt {x} + 2x^{ - \frac {1}{2}} = 5$ can be found by solving: 


 $\textbf{(A)}\ 16x^2-92x+1 = 0\qquad\textbf{(B)}\ 4x^2-25x+4 = 0\qquad\textbf{(C)}\ 4x^2-17x+4 = 0\\ \textbf{(D)}\ 2x^2-21x+2 = 0\qquad\textbf{(E)}\ 4x^2-25x-4 = 0$


	\item (\textit{AHSME 1954}) The expression $\frac{2x^2-x}{(x+1)(x-2)}-\frac{4+x}{(x+1)(x-2)}$ cannot be evaluated for $x=-1$ or $x=2$, 
since division by zero is not allowed. For other values of $x$: 


 $\textbf{(A)}\ \text{The expression takes on many different values.}\\ \textbf{(B)}\ \text{The expression has only the value 2.}\\ \textbf{(C)}\ \text{The expression has only the value 1.}\\ \textbf{(D)}\ \text{The expression always has a value between }-1\text{ and }+2.\\ \textbf{(E)}\ \text{The expression has a value greater than 2 or less than }-1.$


	\item (\textit{AHSME 1954}) If the margin made on an article costing $C$ dollars and selling for $S$ dollars is $M=\frac{1}{n}C$, then the margin is given by: 


 $\textbf{(A)}\ M=\frac{1}{n-1}S\qquad\textbf{(B)}\ M=\frac{1}{n}S\qquad\textbf{(C)}\ M=\frac{n}{n+1}S\\ \textbf{(D)}\ M=\frac{1}{n+1}S\qquad\textbf{(E)}\ M=\frac{n}{n-1}S$


	\item (\textit{AHSME 1954}) The values of $k$ for which the equation $2x^2-kx+x+8=0$ will have real and equal roots are: 


 $\textbf{(A)}\ 9\text{ and }-7\qquad\textbf{(B)}\ \text{only }-7\qquad\textbf{(C)}\ \text{9 and 7}\\ \textbf{(D)}\ -9\text{ and }-7\qquad\textbf{(E)}\ \text{only 9}$


	\item (\textit{AHSME 1954}) The two roots of the equation $a(b-c)x^2+b(c-a)x+c(a-b)=0$ are $1$ and: 


 $\textbf{(A)}\ \frac{b(c-a)}{a(b-c)}\qquad\textbf{(B)}\ \frac{a(b-c)}{c(a-b)}\qquad\textbf{(C)}\ \frac{a(b-c)}{b(c-a)}\qquad\textbf{(D)}\ \frac{c(a-b)}{a(b-c)}\qquad\textbf{(E)}\ \frac{c(a-b)}{b(c-a)}$


	\item (\textit{AHSME 1954}) The straight line $\overline{AB}$ is divided at $C$ so that $AC=3CB$. Circles are described on $\overline{AC}$ 
and $\overline{CB}$ as diameters and a common tangent meets $AB$ produced at $D$. Then $BD$ equals: 


 $\textbf{(A)}\ \text{diameter of the smaller circle}\\ \textbf{(B)}\ \text{radius of the smaller circle}\\ \textbf{(C)}\ \text{radius of the larger circle}\\ \textbf{(D)}\ CB\sqrt{3}\\ \textbf{(E)}\ \text{the difference of the two radii}$


	\item (\textit{AHSME 1954}) A right circular cone has for its base a circle having the same radius as a given sphere. 
The volume of the cone is one-half that of the sphere. The ratio of the altitude of the cone to the radius of its base is:


 $\textbf{(A)}\ \frac{1}{1}\qquad\textbf{(B)}\ \frac{1}{2}\qquad\textbf{(C)}\ \frac{2}{3}\qquad\textbf{(D)}\ \frac{2}{1}\qquad\textbf{(E)}\ \sqrt{\frac{5}{4}}$


	\item (\textit{AHSME 1954}) If $\frac{m}{n}=\frac{4}{3}$ and $\frac{r}{t}=\frac{9}{14}$, the value of $\frac{3mr-nt}{4nt-7mr}$ is: 


 $\textbf{(A)}\ -5\frac{1}{2}\qquad\textbf{(B)}\ -\frac{11}{14}\qquad\textbf{(C)}\ -1\frac{1}{4}\qquad\textbf{(D)}\ \frac{11}{14}\qquad\textbf{(E)}\ -\frac{2}{3}$


	\item (\textit{AHSME 1954}) If the ratio of the legs of a right triangle is $1: 2$, then the ratio of the corresponding segments of 
the hypotenuse made by a perpendicular upon it from the vertex is: 


 $\textbf{(A)}\ 1: 4\qquad\textbf{(B)}\ 1:\sqrt{2}\qquad\textbf{(C)}\ 1: 2\qquad\textbf{(D)}\ 1:\sqrt{5}\qquad\textbf{(E)}\ 1: 5$


	\item (\textit{AHSME 1954}) $A$ and $B$ together can do a job in $2$ days; $B$ and $C$ can do it in four days; and $A$ and $C$ in $2\frac{2}{5}$ days. 
The number of days required for A to do the job alone is: 


 $\textbf{(A)}\ 1 \qquad \textbf{(B)}\ 3 \qquad \textbf{(C)}\ 6 \qquad \textbf{(D)}\ 12 \qquad \textbf{(E)}\ 2.8$


	\item (\textit{AHSME 1954}) In $\triangle ABC$, $AB=AC$, $\angle A=40^\circ$. Point $O$ is within the triangle with $\angle OBC \cong \angle OCA$. 
The number of degrees in $\angle BOC$ is: 


 $\textbf{(A)}\ 110^{\circ} \qquad \textbf{(B)}\ 35^{\circ} \qquad \textbf{(C)}\ 140^{\circ} \qquad \textbf{(D)}\ 55^{\circ} \qquad \textbf{(E)}\ 70^{\circ}$


	\item (\textit{AHSME 1954}) The factors of $x^4+64$ are: 


 $\textbf{(A)}\ (x^2+8)^2\qquad\textbf{(B)}\ (x^2+8)(x^2-8)\qquad\textbf{(C)}\ (x^2+2x+4)(x^2-8x+16)\\ \textbf{(D)}\ (x^2-4x+8)(x^2-4x-8)\qquad\textbf{(E)}\ (x^2-4x+8)(x^2+4x+8)$


	\item (\textit{AHSME 1954}) A bank charges $\textdollar{6}$ for a loan of $\textdollar{120}$. The borrower receives $\textdollar{114}$ and 
repays the loan in $12$ installments of $\textdollar{10}$ a month. The interest rate is approximately: 


 $\textbf{(A)}\ 5 \% \qquad \textbf{(B)}\ 6 \% \qquad \textbf{(C)}\ 7 \% \qquad \textbf{(D)}\ 9\% \qquad \textbf{(E)}\ 15 \%$


	\item (\textit{AHSME 1954}) The fraction $\frac{1}{3}$: 


 $\textbf{(A)}\ \text{equals 0.33333333}\qquad\textbf{(B)}\ \text{is less than 0.33333333 by }\frac{1}{3\cdot 10^8}\\ \textbf{(C)}\ \text{is less than 0.33333333 by }\frac{1}{3\cdot 10^9}\\ \textbf{(D)}\ \text{is greater than 0.33333333 by }\frac{1}{3\cdot 10^8}\\ \textbf{(E)}\ \text{is greater than 0.33333333 by }\frac{1}{3\cdot 10^9}$


	\item (\textit{AHSME 1954}) In the right triangle shown the sum of the distances $BM$ and $MA$ is equal to the sum of the distances $BC$ and $CA$. 
If $MB = x, CB = h$, and $CA = d$, then $x$ equals: 


 
                    \begin{figure}[H]
                        \centering
                        \includegraphics[width=0.5\textwidth]{Problem/AHSME_I_1954_35_Problem_Figure_1.png}
                    \end{figure}
                    
$\textbf{(A)}\ \frac{hd}{2h+d}\qquad\textbf{(B)}\ d-h\qquad\textbf{(C)}\ \frac{1}{2}d\qquad\textbf{(D)}\ h+d-\sqrt{2d}\qquad\textbf{(E)}\ \sqrt{h^2+d^2}-h$


	\item (\textit{AHSME 1954}) A boat has a speed of $15$ mph in still water. In a stream that has a current of $5$ mph it travels a certain 
distance downstream and returns. The ratio of the average speed for the round trip to the speed in still water is: 


 $\textbf{(A)}\ \frac{5}{4}\qquad\textbf{(B)}\ \frac{1}{1}\qquad\textbf{(C)}\ \frac{8}{9}\qquad\textbf{(D)}\ \frac{7}{8}\qquad\textbf{(E)}\ \frac{9}{8}$


	\item (\textit{AHSME 1954}) Given $\triangle PQR$ with $\overline{RS}$ bisecting $\angle R$, $PQ$ extended to $D$ and $\angle n$ a right angle, then: 


 
                    \begin{figure}[H]
                        \centering
                        \includegraphics[width=0.5\textwidth]{Problem/AHSME_I_1954_37_Problem_Figure_1.png}
                    \end{figure}
                    
$\textbf{(A)}\ \angle m = \frac {1}{2}(\angle p - \angle q) \qquad  \textbf{(B)}\ \angle m = \frac {1}{2}(\angle p + \angle q)$
$\textbf{(C)}\ \angle d =\frac{1}{2}(\angle q+\angle p)\qquad \textbf{(D)}\ \angle d =\frac{1}{2}\angle m\qquad \textbf{(E)}\ \text{none of these is correct}$


	\item (\textit{AHSME 1954}) If $\log 2=.3010$ and $\log 3=.4771$, the value of $x$ when $3^{x+3}=135$ is approximately: 


 $\textbf{(A)}\ 5 \qquad \textbf{(B)}\ 1.47 \qquad \textbf{(C)}\ 1.67 \qquad \textbf{(D)}\ 1.78 \qquad \textbf{(E)}\ 1.63$


	\item (\textit{AHSME 1954}) The locus of the midpoint of a line segment that is drawn from a given external point $P$ to a given circle with center $O$ and radius $r$, is: 


 $\textbf{(A)}\ \text{a straight line perpendicular to }\overline{PO}\\ \textbf{(B)}\ \text{a straight line parallel to }\overline{PO}\\ \textbf{(C)}\ \text{a circle with center }P\text{ and radius }r\\ \textbf{(D)}\ \text{a circle with center at the midpoint of }\overline{PO}\text{ and radius }2r\\ \textbf{(E)}\ \text{a circle with center at the midpoint }\overline{PO}\text{ and radius }\frac{1}{2}r$


	\item (\textit{AHSME 1954}) If $\left (a+\frac{1}{a} \right )^2=3$, then $a^3+\frac{1}{a^3}$ equals: 


 $\textbf{(A)}\ \frac{10\sqrt{3}}{3}\qquad\textbf{(B)}\ 3\sqrt{3}\qquad\textbf{(C)}\ 0\qquad\textbf{(D)}\ 7\sqrt{7}\qquad\textbf{(E)}\ 6\sqrt{3}$


	\item (\textit{AHSME 1954}) The sum of all the roots of $4x^3-8x^2-63x-9=0$ is: 


 $\textbf{(A)}\ 8 \qquad \textbf{(B)}\ 2 \qquad \textbf{(C)}\ -8 \qquad \textbf{(D)}\ -2 \qquad \textbf{(E)}\ 0$


	\item (\textit{AHSME 1954}) Consider the graphs of 

 \[(1)\qquad y=x^2-\frac{1}{2}x+2\] 
and 

 \[(2)\qquad y=x^2+\frac{1}{2}x+2\] 
on the same set of axis. 
These parabolas are exactly the same shape. Then: 


 $\textbf{(A)}\ \text{the graphs coincide.}\\ \textbf{(B)}\ \text{the graph of (1) is lower than the graph of (2).}\\ \textbf{(C)}\ \text{the graph of (1) is to the left of the graph of (2).}\\ \textbf{(D)}\ \text{the graph of (1) is to the right of the graph of (2).}\\ \textbf{(E)}\ \text{the graph of (1) is higher than the graph of (2).}$


	\item (\textit{AHSME 1954}) The hypotenuse of a right triangle is $10$ inches and the radius of the inscribed circle is $1$ inch. The perimeter of the triangle in inches is: 


 $\textbf{(A)}\ 15 \qquad \textbf{(B)}\ 22 \qquad \textbf{(C)}\ 24 \qquad \textbf{(D)}\ 26 \qquad \textbf{(E)}\ 30$


	\item (\textit{AHSME 1954}) A man born in the first half of the nineteenth century was $x$ years old in the year $x^2$. He was born in: 


 $\textbf{(A)}\ 1849 \qquad \textbf{(B)}\ 1825 \qquad \textbf{(C)}\ 1812 \qquad \textbf{(D)}\ 1836 \qquad \textbf{(E)}\ 1806$


	\item (\textit{AHSME 1954}) In a rhombus, $ABCD$, line segments are drawn within the rhombus, parallel to diagonal $BD$, 
and terminated in the sides of the rhombus. A graph is drawn showing the length of a segment 
as a function of its distance from vertex $A$. The graph is: 


 $\textbf{(A)}\ \text{A straight line passing through the origin.}\\ \textbf{(B)}\ \text{A straight line cutting across the upper right quadrant.}\\ \textbf{(C)}\ \text{Two line segments forming an upright V.}\\ \textbf{(D)}\ \text{Two line segments forming an inverted V.}\\ \textbf{(E)}\ \text{None of these.}$


	\item (\textit{AHSME 1954}) In the diagram, if points $A, B$ and $C$ are points of tangency, then $x$ equals: 


 
                    \begin{figure}[H]
                        \centering
                        \includegraphics[width=0.5\textwidth]{Problem/AHSME_I_1954_46_Problem_Figure_1.png}
                    \end{figure}
                    
$\textbf{(A)}\ \frac{3}{16}"\qquad\textbf{(B)}\ \frac{1}{8}"\qquad\textbf{(C)}\ \frac{1}{32}"\qquad\textbf{(D)}\ \frac{3}{32}"\qquad\textbf{(E)}\ \frac{1}{16}"$


	\item (\textit{AHSME 1954}) At the midpoint of line segment $AB$ which is $p$ units long, a perpendicular $MR$ is erected with length $q$ units. 
An arc is described from $R$ with a radius equal to $\frac{1}{2}AB$, meeting $AB$ at $T$. Then $AT$ and $TB$ are the roots of: 


 $\textbf{(A)}\ x^2+px+q^2=0\\ \textbf{(B)}\ x^2-px+q^2=0\\ \textbf{(C)}\ x^2+px-q^2=0\\ \textbf{(D)}\ x^2-px-q^2=0\\ \textbf{(E)}\ x^2-px+q=0$


	\item (\textit{AHSME 1954}) A train, an hour after starting, meets with an accident which detains it a half hour, after which it 
proceeds at $\frac{3}{4}$ of its former rate and arrives $3\frac{1}{2}$ hours late. 
Had the accident happened $90$ miles farther along the line, it would have arrived only $3$ hours late. The length of the trip in miles was: 


 $\textbf{(A)}\ 400 \qquad \textbf{(B)}\ 465 \qquad \textbf{(C)}\ 600 \qquad \textbf{(D)}\ 640 \qquad \textbf{(E)}\ 550$


	\item (\textit{AHSME 1954}) The difference of the squares of two odd numbers is always divisible by $8$. If $a>b$, and $2a+1$ and $2b+1$ are the odd numbers, 
to prove the given statement we put the difference of the squares in the form: 


 $\textbf{(A)}\ (2a+1)^2-(2b+1)^2\\ \textbf{(B)}\ 4a^2-4b^2+4a-4b\\ \textbf{(C)}\ 4[a(a+1)-b(b+1)]\\ \textbf{(D)}\ 4(a-b)(a+b+1)\\ \textbf{(E)}\ 4(a^2+a-b^2-b)$


	\item (\textit{AHSME 1954}) The times between $7$ and $8$ o'clock, correct to the nearest minute, when the hands of a clock will form an angle of $84^{\circ}$ are: 


 $\textbf{(A)}\ \text{7: 23 and 7: 53}\qquad \textbf{(B)}\ \text{7: 20 and 7: 50}\qquad \textbf{(C)}\ \text{7: 22 and 7: 53}\\  \textbf{(D)}\ \text{7: 23 and 7: 52}\qquad \textbf{(E)}\ \text{7: 21 and 7: 49}$


\end{enumerate}